\documentclass[11pt,a4paper]{article}
\usepackage[utf8]{inputenc}
\usepackage[spanish]{babel}
\usepackage{geometry}
\usepackage{enumitem}
\usepackage{titlesec}
\usepackage{fancyhdr}
\usepackage{hyperref}
\usepackage{xcolor}

\geometry{margin=2.5cm}
\setlength{\parindent}{0pt}
\setlength{\parskip}{0.6em}

\titleformat{\section}{\large\bfseries}{\thesection.}{0.5em}{}
\pagestyle{fancy}
\fancyhf{}
\rhead{APLICACIÓN Y SERVICIOS WEB}
\lhead{Enunciado del examen}
\cfoot{\thepage}

\definecolor{azul}{RGB}{0,51,102}

\begin{document}

\begin{center}
{\LARGE \textbf{ENUNCIADO DEL EXAMEN}} \\[0.4cm]
{\large Aplicación y Servicios Web} \\[0.2cm]
{\normalsize API REST con FastAPI, base de datos en la nube y menú CRUD} \\[0.8cm]
{\bfseries Fecha máxima de entrega: domingo 8 de marzo}
\end{center}

\vspace{1cm}

\section{Objetivo}

Desarrollar una aplicación que exponga una \textbf{API REST} con \textbf{FastAPI}, utilice \textbf{Uvicorn} como servidor, se conecte a una base de datos PostgreSQL en \textbf{Neon} con un número de entidades entre 6 y 10, y ofrezca un \textbf{menú por consola} para realizar operaciones CRUD sobre los recursos.

\section{Requisitos técnicos}

\begin{enumerate}[leftmargin=*]
  \item \textbf{API REST}
  \begin{itemize}
    \item Implementar una API con \textbf{FastAPI}.
    \item Servir la aplicación con \textbf{Uvicorn} (incluido en la ejecución del proyecto o por comando).
  \end{itemize}

  \item \textbf{Endpoints}
  \begin{itemize}
    \item Definir \textbf{endpoints} para cada recurso (entidad) que permitan:
    \begin{itemize}
      \item Listar todos los registros (GET).
      \item Obtener un registro por identificador (GET por id).
      \item Crear un registro (POST).
      \item Actualizar un registro (PUT).
      \item Eliminar un registro (DELETE).
    \end{itemize}
    \item Los endpoints deben estar organizados (por ejemplo, en módulos o routers por entidad).
  \end{itemize}

  \item \textbf{Base de datos Neon}
  \begin{itemize}
    \item Usar PostgreSQL en \textbf{Neon} (\url{https://neon.tech}) como base de datos en la nube.
    \item Definir entre \textbf{6 y 10 entidades} (tablas) con relaciones coherentes (claves foráneas si aplica).
    \item Utilizar un ORM (por ejemplo SQLAlchemy) para modelos y conexión.
    \item Gestionar la URL de conexión mediante variables de entorno (por ejemplo \texttt{.env}) y no incluir credenciales en el código.
  \end{itemize}

  \item \textbf{Menú y CRUD}
  \begin{itemize}
    \item Implementar un \textbf{menú por consola} que permita al usuario:
    \begin{itemize}
      \item Elegir la entidad o recurso con el que trabajar.
      \item Listar, crear, ver uno, actualizar y eliminar registros (operaciones CRUD).
    \end{itemize}
    \item El menú debe consumir la API (llamadas HTTP a los endpoints), no acceder directamente a la base de datos.
  \end{itemize}

  \item \textbf{Estilo de código y documentación}
  \begin{itemize}
    \item Respetar las normas \textbf{PEP 8} (estilo de código en Python: indentación, longitud de líneas, nombres, espaciado, etc.).
    \item Incluir \textbf{docstrings} en módulos, clases y funciones relevantes (descripción breve del propósito, parámetros y retorno cuando aplique).
  \end{itemize}
\end{enumerate}

\section{Estructura sugerida del proyecto}

\begin{itemize}[leftmargin=*]
  \item \texttt{entities} o \texttt{models}: definición de las tablas/entidades.
  \item \texttt{database}: configuración de conexión a Neon y sesión.
  \item \texttt{schemas}: modelos Pydantic para validación de request/response (opcional pero recomendado).
  \item \texttt{endpoints} o \texttt{routers}: rutas FastAPI por recurso.
  \item \texttt{crud} (cliente): lógica que llama a los endpoints desde el menú (cliente HTTP).
  \item \texttt{main} o script principal: menú por consola que utiliza el CRUD cliente.
  \item Archivo de aplicación FastAPI (por ejemplo \texttt{app.py}) que registre los routers y arranque con Uvicorn.
\end{itemize}

\section{Trabajo con Git y GitHub}

El desarrollo debe realizarse en un repositorio en \textbf{GitHub}. Se evaluará el uso de un flujo de trabajo con ramas y Pull Requests (PR).

\begin{enumerate}[leftmargin=*]
  \item \textbf{Ramas principales}
  \begin{itemize}
    \item \texttt{dev}: desarrollo.
    \item \texttt{qa}: pruebas / calidad.
    \item \texttt{prod}: producción.
  \end{itemize}

  \item \textbf{Flujo de PR por rama}
  \begin{itemize}
    \item Crear ramas de trabajo (features) y integrar mediante \textbf{Pull Requests}:
    \begin{itemize}
      \item \texttt{feat-dev}: rama de desarrollo; se integra en \texttt{dev} vía PR.
      \item \texttt{feat-qa}: integración hacia \texttt{qa} vía PR (desde \texttt{dev} o rama de release).
      \item \texttt{feat-prod}: integración hacia \texttt{prod} vía PR (desde \texttt{qa} o rama de release).
    \end{itemize}
    \item Cada avance hacia \texttt{dev}, \texttt{qa} y \texttt{prod} debe hacerse mediante PR en GitHub (no merges directos en local sin PR).
  \end{itemize}

  \item \textbf{Evaluación individual por commits}
  \begin{itemize}
    \item La nota en esta parte será \textbf{individual} según el historial de \textbf{commits} en el repositorio.
    \item Se tendrá en cuenta la autoría de los commits, la claridad de los mensajes y la trazabilidad del trabajo (qué aportó cada integrante en las ramas y PR).
  \end{itemize}
\end{enumerate}

\section{Criterios de evaluación}

\begin{itemize}[leftmargin=*]
  \item Funcionamiento de la API con FastAPI y Uvicorn.
  \item Correcta definición y uso de endpoints (verbos HTTP y recursos).
  \item Conexión a Neon y definición de entre 6 y 10 entidades con sentido lógico.
  \item Menú operativo que realice CRUD mediante llamadas a la API.
  \item Organización del código y buenas prácticas (variables de entorno, sin credenciales en código).
  \item Cumplimiento de \textbf{PEP 8} y uso de \textbf{docstrings} en el código.
  \item \textbf{Git/GitHub}: uso de ramas \texttt{dev}, \texttt{qa}, \texttt{prod} y flujo de PR (\texttt{feat-dev}, \texttt{feat-qa}, \texttt{feat-prod}); nota individual según commits.
\end{itemize}

\vspace{0.5cm}
\hrule
\vspace{0.3cm}
{\small Documento generado para el curso de Aplicación y Servicios Web.}

\end{document}
